% Supplementary material for main_final.tex.

\documentclass{vgtc}

\graphicspath{{figures/}{pictures/}{images/}{./}}

\usepackage{times}
\renewcommand*\ttdefault{txtt}
\usepackage{mathptmx}
\usepackage{amsmath}
\usepackage{amssymb}
\usepackage{array}
\usepackage{booktabs}
\usepackage{multirow}
\usepackage{graphicx}
\usepackage{xcolor}
\usepackage{url}
\usepackage{enumitem}
\usepackage{xspace}

\newcommand{\R}{\mathbb{R}}
\newcommand{\Domain}{\mathcal{M}}
\newcommand{\Range}{\mathcal{R}}
\newcommand{\Sheets}{\mathcal{S}}
\newcommand{\Sheet}{S}
\newcommand{\Score}{C}
\newcommand{\placeholder}[2]{%
  \begin{center}
  \fbox{%
    \begin{minipage}[c][#1][c]{0.92\linewidth}
      \centering
      #2
    \end{minipage}}
  \end{center}
}
\newcommand{\wideplaceholder}[2]{%
  \begin{center}
  \fbox{%
    \begin{minipage}[c][#1][c]{0.96\textwidth}
      \centering
      #2
    \end{minipage}}
  \end{center}
}
\newcommand{\myparagraph}[1]{\vspace{1mm} \noindent \textbf{#1}}
\newcommand{\ie}{i.e.,\xspace}
\newcommand{\etal}{et al.\xspace}

\onlineid{0}
\vgtccategory{Supplementary Material}
\vgtcinsertpkg

\title{Supplementary Material: Topological Tracking of Bivariate Features Using Reeb Spaces}

\author{Mohit Sharma\\ %
        \scriptsize Link\"oping university}

\begin{document}

\firstsection{Supplementary Methods}
\maketitle

This supplement gives the implementation-level details behind the main paper. It is intentionally more exhaustive than the main text. The main paper focuses on the visualization workflow and case-study interpretation; this document records the exact score definitions, threshold choices, event diagnostics, disagreement score, and viewer controls used by the current implementation.

\section{Pipeline Stages and Artifacts}

The input is a time-ordered directory of \texttt{.vtu} files. Each file stores two scalar arrays, configured as \texttt{orb00} and \texttt{orb01} in the current experiments. The active implementation is a staged Python pipeline around an external Reeb-space executable and a generated browser viewer. The stages are intentionally separated so that expensive Reeb-space computation, geometry export, image rendering, and paper analysis can be rerun independently.

\myparagraph{Stage 1: Reeb-space computation.}
The first stage scans the input \texttt{.vtu} directory and runs the Reeb-space executable for each timestep. It produces two primary files per successful timestep:
\begin{itemize}[leftmargin=*]
  \item \texttt{.rs}: Reeb-space geometry and sheet structure used for range-space geometry export and rendering.
  \item \texttt{.rsi}: sheet information used by the Python stages, including sheet identifiers, sheet areas, and associated regular domain vertices.
\end{itemize}
Only timesteps with usable Reeb-space sheet information are included in the temporal graph. Failed or missing timesteps are logged and skipped by downstream stages rather than being filled with synthetic matches.

\myparagraph{Stage 2: RSI-to-JSON conversion.}
The second stage converts each binary \texttt{.rsi} file into a compact \texttt{.rsijson} summary. For each selected sheet, the JSON stores the sheet id, rank, range-space area, number of stored regular domain vertices, and the regular vertex set. This is the lightweight representation used by domain-overlap computation and viewer metadata.

\myparagraph{Stage 3A: range-shape matching.}
The range matching stage exports sheet geometry from the Reeb-space files and caches reusable VTP and mask data. It compares selected source sheets and target sheets by rasterizing their range-space footprints under common global bounds. The main output is \texttt{sheet\_shape\_matches.json}, with per-link range metrics including shape IoU, area ratio, bounding-box IoU, centroid similarity, and the current combined range score.

\myparagraph{Stage 3B: domain-overlap matching.}
The domain-overlap stage reads the \texttt{.rsijson} files and computes vertex-set intersections between selected source and target sheets. For a source sheet and a target sheet, it stores the raw shared-vertex count and directional overlap ratios. The output \texttt{sheet\_overlaps.json} also attaches available range metrics for matching sheet pairs, which lets the viewer compare domain and range evidence for the same candidate correspondence.

\myparagraph{Stage 4A: sheet rendering.}
The sheet-rendering stage creates range-space sheet images. It reuses cached VTP geometry and renders all images with one global range-space extent and a fixed image size. This is important: if every sheet were rendered with its own camera, two images could appear similarly sized even when the underlying range footprints differ strongly.

\myparagraph{Stage 4B: fiber-surface rendering.}
The fiber-surface stage renders sheet-associated fiber-surface images using the configured fiber-surface isovalues and rendering state. These images provide spatial context for selected sheets and links. Empty or missing surfaces are preserved as missing evidence; they are not interpreted as successful continuations.

\myparagraph{Stage 5: unified viewer generation.}
The final viewer stage writes a self-contained browser directory containing \texttt{index.html}, \texttt{style.css}, \texttt{viewer.js}, \texttt{viewer\_common.js}, \texttt{data.json}, and links or copies to sheet and fiber-surface image directories. The browser entry point is the generated \texttt{unified\_sankey\_viewer} directory.

\myparagraph{Tracking-analysis stage.}
The analysis stage consumes the generated viewer data and writes event, lifetime, sensitivity, metric-agreement, and domain/range disagreement summaries. It does not recompute Reeb spaces, VTP geometry, sheet images, or fiber-surface images. This separation keeps the reported analysis tied to a fixed generated viewer dataset.

\section{Sheet Representation and Top-$N$ Selection}

For timestep $i$, let $\Sheets_i^\mathrm{all}$ be all computed Reeb-space sheets. Each sheet stores a sheet id, range-space area, range-space polygonal footprint, rank, regular domain-vertex set, centroid, and bounding box. The viewer keeps the top $N$ sheets by range-space area:
$$
\Sheets_i=\operatorname{TopN}(\Sheets_i^\mathrm{all},N).
$$
The current experiments use $N=20$. The selected set is recomputed independently at each timestep, so a sheet id may appear, disappear, or change rank over time. Sheets with zero stored regular vertices are retained. Such sheets may still have nonzero range-space area, but all domain-overlap scores involving their empty domain-vertex set are zero.

\section{Range-Space Correspondence Scores}

Range-space scores compare sheet footprints in the bivariate range. For a pair of sheets $\Sheet\in\Sheets_i$ and $\Sheet'\in\Sheets_{i+s}$, where $s$ is the timestep stride, the implementation rasterizes each sheet footprint into a mask over common global range bounds. The same global bounds and grid size are used for all masks in a dataset, so mask overlap is comparable across timesteps.

\myparagraph{Shape IoU.}
Let $M_\Sheet$ and $M_{\Sheet'}$ be the two binary masks. The shape intersection-over-union score is
$$
q_\mathrm{iou}(\Sheet,\Sheet')=
\frac{|M_\Sheet\cap M_{\Sheet'}|}{|M_\Sheet\cup M_{\Sheet'}|}.
$$
This is the current primary range-space score. A high value means that the two Reeb-space sheets occupy similar regions of the bivariate range. Shape IoU is preferred for the main analysis because it directly measures detailed footprint overlap under a common coordinate frame. It already penalizes shifts, large size changes, and detailed shape deformation through the intersection and union terms.

\myparagraph{Area ratio.}
Let $A_\Sheet$ and $A_{\Sheet'}$ be the range-space areas. The area-ratio score is
$$
q_\mathrm{area}(\Sheet,\Sheet')=
\frac{\min(A_\Sheet,A_{\Sheet'})}{\max(A_\Sheet,A_{\Sheet'})},
$$
with score zero if both areas are invalid or nonpositive. This score checks size consistency only; it does not check position or shape.

\myparagraph{Bounding-box IoU.}
Let $B_\Sheet$ and $B_{\Sheet'}$ be axis-aligned bounding boxes in the range plane. The bounding-box score is
$$
q_\mathrm{bbox}(\Sheet,\Sheet')=
\frac{\operatorname{area}(B_\Sheet\cap B_{\Sheet'})}{\operatorname{area}(B_\Sheet\cup B_{\Sheet'})}.
$$
It is cheaper and coarser than shape IoU. It can agree with shape IoU when sheets are compact and similarly shaped, but it can also be misleading for nonconvex or elongated footprints.

\myparagraph{Centroid similarity.}
Let $c_\Sheet$ and $c_{\Sheet'}$ be range-space centroids and let $D$ be the diagonal length of the global range-space bounds. The centroid score is
$$
q_\mathrm{centroid}(\Sheet,\Sheet')=
\exp\left(-\frac{\|c_\Sheet-c_{\Sheet'}\|}{D}\right).
$$
This score decreases smoothly with centroid distance. It is useful as a weak positional cue, but it is not discriminative enough to be the primary matching score.

\myparagraph{Combined range score.}
The code supports a weighted combined range score:
$$
q_R(\Sheet,\Sheet')=
\frac{
w_\mathrm{iou}q_\mathrm{iou}
+w_\mathrm{area}q_\mathrm{area}
+w_\mathrm{bbox}q_\mathrm{bbox}
+w_\mathrm{centroid}q_\mathrm{centroid}
}{
w_\mathrm{iou}+w_\mathrm{area}+w_\mathrm{bbox}+w_\mathrm{centroid}
}.
$$
In the current generated results, $w_\mathrm{iou}=1$ and the other weights are zero. Therefore $q_R=q_\mathrm{iou}$ in the case studies. The auxiliary metrics remain visible in the viewer and in metric diagnostics, but they are not used to change the main range-event ranking. This choice avoids mixing weaker cues into the primary event score: area ratio ignores location, centroid similarity ignores footprint shape, and bounding-box IoU is a coarse approximation of the actual sheet mask.

\section{Domain-Space Correspondence Scores}

Domain scores compare regular domain-vertex sets. Let $V_\Sheet$ and $V_{\Sheet'}$ be the regular vertices associated with source and target sheets. The raw overlap count is
$$
O(\Sheet,\Sheet')=|V_\Sheet\cap V_{\Sheet'}|.
$$
This value asks whether the same domain vertices continue to participate in the sheet behavior.

\myparagraph{Why not vertex-set Jaccard overlap?}
A natural symmetric score is
$$
q_\mathrm{jac}(\Sheet,\Sheet')=
\frac{|V_\Sheet\cap V_{\Sheet'}|}{|V_\Sheet\cup V_{\Sheet'}|}.
$$
This score answers a strict equality-style question: are the two domain-vertex sets nearly the same? That is not always the right question for Reeb-space sheet tracking. Suppose all vertices of a small source sheet continue into a larger target sheet. Then the source-retention evidence is perfect, but Jaccard can be small because the union includes many additional target vertices. The same problem appears in the opposite direction when a large source shrinks into a smaller target, and in partial merge or absorption events where containment is meaningful evidence rather than a failure. For this reason, Jaccard overlap is useful as a reference point but is not used as the reported domain event or disagreement score.

The implementation stores directional overlap ratios:
$$
d_\mathrm{src}(\Sheet,\Sheet')=
\begin{cases}
O(\Sheet,\Sheet')/|V_\Sheet|, & |V_\Sheet|>0,\\
0, & |V_\Sheet|=0,
\end{cases}
$$
and
$$
d_\mathrm{tgt}(\Sheet,\Sheet')=
\begin{cases}
O(\Sheet,\Sheet')/|V_{\Sheet'}|, & |V_{\Sheet'}|>0,\\
0, & |V_{\Sheet'}|=0.
\end{cases}
$$
The domain-overlap metric used in the event and disagreement views is
$$
d_\mathrm{max}(\Sheet,\Sheet')=\max(d_\mathrm{src},d_\mathrm{tgt}).
$$
This score intentionally treats containment as strong continuity. If a small source sheet is contained in a larger target sheet, $d_\mathrm{src}$ can be high even if the target has additional vertices. This is useful for growth, shrinkage, absorption, and partial merge/split behavior. We use
$$
q_D(\Sheet,\Sheet')=d_\mathrm{max}(\Sheet,\Sheet'),
$$
so that domain and range scores both lie in $[0,1]$. The interface may format these ratios as percentages for readability, but the analysis formulas use ratios.

\section{Thresholds and Percentiles}

\myparagraph{Fixed range thresholds.}
Range event and lifetime summaries are exported for thresholds
$$
\theta\in\{0.3,0.4,0.5,0.6,0.7\}.
$$
The main paper reports $\theta=0.5$ unless stated otherwise. For shape IoU, this means that the intersection of the two range-footprint masks is at least half as large as their union. It is therefore a mid-scale exploratory cutoff for substantial range-footprint continuity, not a physically calibrated constant.

\myparagraph{Runtime domain thresholds.}
Domain scores depend on the mesh, domain-vertex set sizes, and selected domain metric. For interactive domain views, the viewer derives threshold options from the observed best domain continuations. Let $q_D^\mathrm{act}$ be the active domain score converted to $[0,1]$. For source, target, and maximum directional domain ratios, this is the corresponding ratio. For raw shared-vertex count, this is the count normalized by the maximum shared-vertex count for the active dataset and metric. For each source sheet $\Sheet\in\Sheets_i$, the viewer computes
$$
b_D(\Sheet)=\max_{\Sheet'\in\Sheets_{i+s}}q_D^\mathrm{act}(\Sheet,\Sheet'),
$$
where $s$ is the selected timestep stride. The viewer collects all such values into a list
$$
\mathcal{Q}_D=\{b_D(\Sheet):\Sheet\in\Sheets_i,\ \mathrm{all\ visible\ source\ timesteps}\}.
$$
After removing non-finite values, the list is sorted:
$$
z_0\leq z_1\leq\cdots\leq z_{m-1}.
$$
For a fraction $\alpha$, the implementation uses nearest-rank indexing with rounding:
$$
Q_\alpha=z_{\operatorname{round}(\alpha(m-1))}.
$$
The list $\mathcal{Q}_D$ should be read as the empirical distribution of how good each source sheet's best available domain continuation is. If a dataset has generally weak domain overlap, a fixed threshold such as $0.5$ may remove nearly all domain links. If a dataset has generally strong domain overlap, the same fixed threshold may show too many links. Percentiles adapt the threshold to the observed distribution: $Q_{50}$ keeps links that are at least as strong as a typical best continuation, while $Q_{75}$ and $Q_{90}$ show progressively stronger and more conservative domain continuations. The viewer offers $Q_{25}$, $Q_{50}$, $Q_{75}$, and $Q_{90}$ as domain-threshold options, and uses $Q_{50}$ as the default.

\section{Temporal Graph}

For each stride $s$, the temporal graph compares sheets between timestep $i$ and timestep $i+s$. For a fixed pair of timesteps, the graph is a bipartite candidate graph:
$$
G_i^s=(\Sheets_i,\Sheets_{i+s},E_i^s).
$$
Every edge stores range metrics, domain metrics, and metadata for source and target sheets. The viewer displays the candidate graph in either shape-overlap mode or domain-overlap mode. Keeping these modes separate is deliberate: range continuity asks whether the bivariate field relationship persists, while domain overlap asks whether the same regular vertices continue to participate in the sheet.

No global one-to-one assignment is solved. Multiple source sheets can select the same target sheet, and one source can have several plausible targets. This is important because Reeb-space sheets may split, merge, appear, disappear, or change shape. The Sankey graph is therefore a visual summary of candidate correspondence evidence, not a final identity assignment.

\section{Sankey Encoding and Interaction}

The generated viewer has three main regions: a control panel, one or more Sankey panels, and a detail panel. The control panel manages timestep ranges, layout, top-sheet count, timestep stride, node coloring, link darkness, hidden isolated nodes, strongest-outgoing filtering, and figure-preset export. The center region contains the range bar and the Sankey panels. The detail panel reports the metadata and visual evidence for the currently selected node or link.

\myparagraph{Timestep ranges and stride.}
The range bar lets the analyst focus on one or more timestep intervals. The user can add ranges, select a range, drag range endpoints, delete a range, and center the current view. The timestep stride control changes the temporal sampling used by the visible graph. A stride of $s$ compares timestep $i$ to timestep $i+s$ where precomputed stride data are available. This is used in the results to check whether important intervals remain stable under coarser sampling.

\myparagraph{Panels and active metric.}
Each Sankey panel has its own data mode, metric, threshold, and height. Shape-overlap mode can display shape IoU, the current combined shape score, area ratio, bounding-box IoU, or centroid similarity. Domain-overlap mode can display shared domain-vertex count, source-domain ratio, target-domain ratio, or maximum domain ratio. Multiple panels can be shown together, which makes it possible to compare shape-overlap and domain-overlap behavior over the same selected timestep window.

\myparagraph{Node height.}
Node height encodes the number of stored regular vertices associated with a sheet, normalized within the visible layout. Zero-vertex sheets are drawn with a minimum height so they remain visible and selectable. Since vertical space is limited, node height should be read as a visual cue rather than as an exact bar chart.

\myparagraph{Link width and opacity.}
For the active metric $q$, a visible link is normalized by the maximum active score in the current data:
$$
\widehat{q}(\Sheet,\Sheet')=
\frac{q(\Sheet,\Sheet')}{\max_{\Sheet_a,\Sheet_b}q(\Sheet_a,\Sheet_b)}.
$$
The normalized score is clamped to $[0,1]$ and used for link width and opacity. Stronger correspondences are thicker and more opaque.

\myparagraph{Threshold filtering.}
The link threshold is applied to the normalized active score:
$$
\widehat{q}(\Sheet,\Sheet')\geq\theta.
$$
Changing the active metric changes the meaning of the threshold. A range threshold filters by range-shape continuity, while a domain threshold filters by domain-overlap continuity.

\myparagraph{Strongest outgoing links.}
The viewer can show only the strongest outgoing link from each source sheet under the active metric. This reduces clutter but does not impose a one-to-one matching. Several source sheets may still select the same target sheet.

\myparagraph{Node ordering and coloring.}
Sheets can be ordered by weighted crossing-reduced ordering, range-space area, or domain vertex count. The crossing-reduced ordering uses the visible links and their visual weights to reduce link crossings while preserving the left-to-right timestep structure. Color can encode a neutral sheet color, area, vertex count, or centroid position. Centroid coloring uses global range bounds so that colors are comparable across time. The current viewer includes a four-corner centroid color map and an axis-oriented color map where saturation increases with distance from the origin.

\myparagraph{Linked detail views.}
Selecting a node shows sheet metadata, the sheet footprint image, and available fiber-surface images. Selecting a link shows source and target sheets side by side. The image viewer supports linked pan and zoom for comparing corresponding sheets and fiber surfaces.

\myparagraph{Analysis tabs.}
The analysis box is available inside each panel. In shape-overlap mode, the interval tab reports interesting range intervals; in domain-overlap mode, it reports interesting domain intervals. The continuing-feature tab lists long greedy tracks above the selected threshold. The sensitivity tab summarizes how top intervals and lifetimes change across thresholds. The domain/range complementarity tab is available from either mode and highlights source sheets for which the best shape-overlap target differs from the best domain-overlap target. Highlight buttons expand the visible ranges when needed and mark the corresponding nodes and links in the Sankey graph.

\section{Event Score Details}

Let $q$ be the active score and let $N_i=|\Sheets_i|$. For a source sheet $\Sheet\in\Sheets_i$, define the best outgoing score
$$
B^\mathrm{out}(\Sheet)=\max_{\Sheet'\in\Sheets_{i+s}}q(\Sheet,\Sheet').
$$
For a target sheet $\Sheet'\in\Sheets_{i+s}$, define the best incoming score
$$
B^\mathrm{in}(\Sheet')=\max_{\Sheet\in\Sheets_i}q(\Sheet,\Sheet').
$$
At threshold $\theta$, the weak outgoing and weak incoming counts are
$$
W^\mathrm{out}_i(\theta)=
\left|\{\Sheet\in\Sheets_i:B^\mathrm{out}(\Sheet)<\theta\}\right|
$$
and
$$
W^\mathrm{in}_i(\theta)=
\left|\{\Sheet'\in\Sheets_{i+s}:B^\mathrm{in}(\Sheet')<\theta\}\right|.
$$

A possible split is counted when one source sheet has at least two target sheets above threshold:
$$
P^\mathrm{split}_i(\theta)=
\left|
\left\{
\Sheet\in\Sheets_i:
\left|\{\Sheet'\in\Sheets_{i+s}:q(\Sheet,\Sheet')\geq\theta\}\right|\geq2
\right\}
\right|.
$$
A possible merge is counted when one target sheet has at least two source sheets above threshold:
$$
P^\mathrm{merge}_i(\theta)=
\left|
\left\{
\Sheet'\in\Sheets_{i+s}:
\left|\{\Sheet\in\Sheets_i:q(\Sheet,\Sheet')\geq\theta\}\right|\geq2
\right\}
\right|.
$$

The mean best outgoing score is
$$
\overline{B^\mathrm{out}_i}=
\frac{1}{N_i}\sum_{\Sheet\in\Sheets_i}B^\mathrm{out}(\Sheet).
$$
The event score is
$$
E_i(\theta)=
W^\mathrm{out}_i(\theta)+W^\mathrm{in}_i(\theta)
+0.5P^\mathrm{split}_i(\theta)+0.5P^\mathrm{merge}_i(\theta)
+(1-\overline{B^\mathrm{out}_i})N_i.
$$
This is a ranking score, not a probability. The first two terms emphasize missing or weak continuations. The split and merge terms emphasize branching ambiguity, but with half weight because a sheet involved in a split or merge still has above-threshold matches. The final continuation-gap term catches broad weakening even if most sheets barely pass the threshold. A high event score therefore means that the interval deserves inspection because continuity is poor, ambiguous, or both.

Range intervals in the paper use $q=q_R$ and answer: did the range-space sheet footprints change? Domain intervals reported in the paper use $q=q_D$, the maximum directional domain-overlap ratio, and answer: did the overlap among regular domain vertices change? In the viewer, the same interval computation can also be run with the active domain metric, such as raw shared-vertex count, source-domain ratio, or target-domain ratio. These questions are intentionally different from range matching.

\section{Continuing-Feature Lifetimes}

For a sheet $\Sheet_i\in\Sheets_i$, the greedy lifetime procedure selects
$$
\Sheet_{i+s}^\ast=\arg\max_{\Sheet'\in\Sheets_{i+s}}q(\Sheet_i,\Sheet').
$$
If $q(\Sheet_i,\Sheet_{i+s}^\ast)\geq\theta$, the track continues to $\Sheet_{i+s}^\ast$ and the same rule is applied again. The lifetime is the number of visited timesteps before the best continuation falls below threshold or the sequence ends.

The track records length, rank range, mean continuation score, and minimum continuation score. It is a diagnostic for persistent behavior, not a proof of unique identity. Multiple starting sheets may converge to the same later sheet.

\section{Domain/Range Target Disagreement}

For each source sheet $\Sheet\in\Sheets_i$, the viewer computes the best range target
$$
T_R(\Sheet)=\arg\max_{\Sheet'\in\Sheets_{i+s}}q_R(\Sheet,\Sheet')
$$
and the best domain target
$$
T_D(\Sheet)=\arg\max_{\Sheet'\in\Sheets_{i+s}}q_D(\Sheet,\Sheet').
$$
If $T_R(\Sheet)=T_D(\Sheet)$, range and domain agree for that source sheet. If the targets differ, the source sheet contributes a disagreement example.

Let
$$
q_R^\ast=q_R(\Sheet,T_R(\Sheet)),
\qquad
q_D^\ast=q_D(\Sheet,T_D(\Sheet)).
$$
The cross alternatives are
$$
q_R^D=q_R(\Sheet,T_D(\Sheet)),
\qquad
q_D^R=q_D(\Sheet,T_R(\Sheet)).
$$
The losses are
$$
L_R=\max(0,q_R^\ast-q_R^D),
\qquad
L_D=\max(0,q_D^\ast-q_D^R).
$$
The confidence term is
$$
C=\min(q_R^\ast,q_D^\ast).
$$
The source-level disagreement score is
$$
D(\Sheet)=0.5(L_R+L_D)C.
$$

This score is high only when both modes have confident preferred targets and each mode loses score by accepting the other mode's target. It is low when the two targets are almost interchangeable or when one mode is already uncertain.

For a timestep pair, the viewer summarizes all source-level disagreements by reporting the disagreement count, agreement fraction, mean disagreement score, and maximum disagreement score:
$$
D_\mathrm{max}(i,i+s)=\max_{\Sheet\in\Sheets_i}D(\Sheet).
$$
The y-axis in the domain/range disagreement graph is this maximum value. Therefore, a high point means that at least one source sheet in that interval has a strong local disagreement between range and domain targets. It does not mean that every sheet in the interval disagrees, and it should be interpreted together with the disagreement count and agreement fraction.

\section{Metric Agreement Diagnostics}

The analysis compares the best target selected by each candidate metric against the best target selected by the current combined range score. Because the current combined score equals shape IoU, this diagnostic asks which metrics choose the same target as shape IoU and how much score would be lost if a different metric were used.

\begin{table*}[t]
  \centering
  \caption{Best-target agreement with the current combined range score. Domain overlap is intentionally low because it asks a different question from range-shape continuity.}
  \label{tab:supp-metric-agreement}
  \begin{tabular}{llrr}
    \toprule
    Dataset & Candidate metric & Agreement & Mean loss if used \\
    \midrule
    \texttt{MVK\_s1} & shape IoU & 1.000 & 0.000 \\
    \texttt{MVK\_s1} & bbox IoU & 0.951 & 0.016 \\
    \texttt{MVK\_s1} & centroid similarity & 0.921 & 0.028 \\
    \texttt{MVK\_s1} & area ratio & 0.815 & 0.098 \\
    \texttt{MVK\_s1} & domain max overlap & 0.112 & 0.630 \\
    \texttt{MVK\_s2} & shape IoU & 1.000 & 0.000 \\
    \texttt{MVK\_s2} & bbox IoU & 0.931 & 0.014 \\
    \texttt{MVK\_s2} & centroid similarity & 0.897 & 0.029 \\
    \texttt{MVK\_s2} & area ratio & 0.733 & 0.134 \\
    \texttt{MVK\_s2} & domain max overlap & 0.077 & 0.627 \\
    \texttt{stilbene} & shape IoU & 1.000 & 0.000 \\
    \texttt{stilbene} & bbox IoU & 0.882 & 0.032 \\
    \texttt{stilbene} & centroid similarity & 0.849 & 0.043 \\
    \texttt{stilbene} & area ratio & 0.568 & 0.232 \\
    \texttt{stilbene} & domain max overlap & 0.062 & 0.589 \\
    \texttt{torus} & shape IoU & 1.000 & 0.000 \\
    \texttt{torus} & bbox IoU & 0.598 & 0.000 \\
    \texttt{torus} & centroid similarity & 0.562 & 0.016 \\
    \texttt{torus} & area ratio & 0.535 & 0.013 \\
    \texttt{torus} & domain max overlap & 0.175 & 0.265 \\
    \bottomrule
  \end{tabular}
\end{table*}

The metric agreement table supports the main-paper interpretation. Shape IoU is the primary range metric. Bounding-box IoU and centroid similarity are useful secondary diagnostics in some datasets. Area ratio is weaker because it ignores position. Domain overlap should not be treated as a substitute for range matching because it measures persistence of overlapping regular domain vertices rather than persistence of range-space footprint.

\section{Recommended Figure Captures}

For reproducible figure generation, the following viewer captures are recommended:
\begin{itemize}[leftmargin=*]
  \item MVK range-event overview: range mode, shape IoU, $\theta=0.5$, intervals around $1180$--$1280$ for \texttt{MVK\_s1} and $1140$--$1180$ for \texttt{MVK\_s2}.
  \item MVK disagreement detail: select a high disagreement interval near $1240$--$1248$ and show the source sheet, range-selected target, domain-selected target, and fiber-surface images.
  \item Stilbene range-event overview: show $9380$--$9680$ and $11900$--$11935$ windows with event graph and highlighted Sankey columns.
  \item Stilbene long-lived tracks: highlight tracks $86\rightarrow6270$, $38\rightarrow19239$, and $19\rightarrow53844$.
  \item Torus validation: show the event graph over all timesteps with a zoom around $45$--$55$ and stride annotations centered at timestep 50.
  \item Torus Sankey detail: show the range-mode Sankey around $45$--$55$ and one full-length track such as $379\rightarrow43$.
\end{itemize}

\section{Known Interpretation Caveats}

The event score is an exploratory diagnostic. It prioritizes intervals for inspection but is not a formal theorem about topological events. The domain score is based on regular vertices stored in the RSI representation; it does not fully capture singular-structure membership. Zero-vertex sheets can be meaningful range-space features but provide no regular-vertex domain evidence. Finally, adjacent-pair tracks can fragment or converge because the method intentionally avoids a global one-to-one optimization.

\bibliographystyle{abbrv-doi}
\bibliography{bib}

\end{document}
